\clearpage\pagestyle{otMain}\setlength{\otBlocEspacement}{2.5mm}\fontsize{11}{15}\selectfont\corrigesChapitre{10}{Longueurs et périmètres}
\begin{blocLivre}{cor-10-1}\begin{corrige}{1}{Convertir en mètres}\demarche{Multiplier par la valeur de l’unité en mètres.}
$\res{3\,700}$ m ; $\res{0{,}45}$ m ; $\res{200}$ m ; $\res{0{,}008}$ m ; $\res{60}$ m.
\end{corrige}\end{blocLivre}
\begin{blocLivre}{cor-10-2}\begin{corrige}{2}{Convertir depuis le mètre}\demarche{1 m vaut 100 cm, 1 000 mm ; 1 km vaut 1 000 m.}
$\res{80}$ cm ; $\res{1{,}25}$ km ; $\res{2\,500}$ mm ; $\res3$ dam.
\end{corrige}\end{blocLivre}
\begin{blocLivre}{cor-10-3}\begin{corrige}{3}{Carré et rectangle}\demarche{Remplacer les lettres des formules par les mesures.}
Carré : $4\times3{,}5=\res{14\ \mathrm{cm}}$. Rectangle : $2\times(7{,}2+4{,}5)=2\times11{,}7=\res{23{,}4\ \mathrm{m}}$.
\end{corrige}\end{blocLivre}
\begin{blocLivre}{cor-10-4}\begin{corrige}{4}{Périmètre d’un disque}\demarche{Le diamètre vaut deux rayons.}
$\pi\times10$ cm, soit environ $\res{31{,}4\ \mathrm{cm}}$. Rayon 5 m : diamètre 10 m, périmètre $10\pi$ m, environ $\res{31{,}4\ \mathrm{m}}$.
\end{corrige}\end{blocLivre}
\begin{blocLivre}{cor-10-5}\begin{corrige}{5}{Reporter au compas}\demarche{Le compas transporte une longueur sans la mesurer.}
Les reports figurent en vert. On obtient $2{,}3+1{,}8+2{,}6=\res{6{,}7\ \mathrm{cm}}$ : c’est la longueur de la ligne brisée.
\label{fig-cor-10-5}\figureGeo[corrige]{lon-5}
\end{corrige}\end{blocLivre}
\begin{blocLivre}{cor-10-6}\begin{corrige}{6}{Une figure composée}\demarche{Ne pas compter le segment intérieur.}
Côtés droits : $6+4+6=16$ cm. Demi-cercle : $3{,}14\times4:2=6{,}28$ cm. Périmètre : environ $\res{22{,}28\ \mathrm{cm}}$.
\label{fig-cor-10-6}\figureGeo[corrige]{lon-6}
\end{corrige}\end{blocLivre}
\begin{blocLivre}{cor-10-7}\begin{corrige}{7}{Le tour du stade}\demarche{Assembler les deux demi-cercles.}
Les deux demi-cercles forment un cercle entier : $3{,}14\times60=188{,}4$ m. Un tour : $2\times100+188{,}4=\res{388{,}4\ \mathrm{m}}$ environ.
\label{fig-cor-10-7}\figureGeo[corrige]{lon-7}
\end{corrige}\end{blocLivre}
\begin{blocLivre}{cor-10-8}\begin{corrige}{8}{Retrouver une dimension}\demarche{Le demi-périmètre vaut L plus l.}
Une longueur plus une largeur font $30:2=15$ cm : largeur $15-9=\res6$ cm. Côté du carré : $26:4=\res{6{,}5}$ cm.
\end{corrige}\end{blocLivre}
\begin{blocLivre}{cor-10-9}\begin{corrige}{9}{Proportionnel au diamètre}\demarche{Multiplier chaque diamètre par 3,14.}
$\res{3{,}14}$ ; $\res{6{,}28}$ ; $\res{15{,}7}$ ; $\res{31{,}4}$ ; $\res{62{,}8}$. Quand le diamètre est multiplié par 3, le périmètre est aussi multiplié par 3 : c’est une situation de proportionnalité.
\end{corrige}\end{blocLivre}
\begin{blocLivre}{cor-10-10}\begin{corrige}{10}{Le contour d’un L}\demarche{Utiliser les côtés opposés de même direction.}
Côté horizontal manquant : $8-3=\res5$ cm. Côté vertical manquant : $7-3=\res4$ cm. Périmètre : $8+3+5+4+3+7=\res{30\ \mathrm{cm}}$.
\label{fig-cor-10-10}\figureGeo[corrige]{lon-10}
\end{corrige}\end{blocLivre}
\begin{blocLivre}{cor-10-11}\begin{corrige}{11}{La roue du vélo}\demarche{Un tour de roue fait avancer d’un périmètre.}
Un tour : $3{,}14\times70=219{,}8$ cm, soit $\res{2{,}198\ \mathrm{m}}$ environ. 1 000 tours : $2\,198$ m, soit environ $\res{2{,}2\ \mathrm{km}}$.
\end{corrige}\end{blocLivre}
\begin{blocLivre}{cor-10-12}\begin{corrige}{12}{La clôture}\demarche{Retirer la longueur du portail.}
Périmètre : $2\times(18{,}5+12)=61$ m. Grillage : $61-1{,}5=\res{59{,}5\ \mathrm{m}}$. Prix : $59{,}5\times2{,}40=\res{142{,}80}$ €.
\end{corrige}\end{blocLivre}
\begin{blocLivre}{cor-10-13}\begin{corrige}{13}{Le marathon}\demarche{Convertir, puis effectuer une division euclidienne.}
$42{,}195$ km $=42\,195$ m et $42\,195=400\times105+195$ : $\res{105}$ tours complets et $\res{195\ \mathrm{m}}$ de plus.
\end{corrige}\end{blocLivre}
\begin{blocLivre}{cor-10-14}\begin{corrige}{14}{Le rond-point}\demarche{Calculer le tour entier, puis la fraction.}
Tour complet : $2\times3{,}14\times15=94{,}2$ m. Trois quarts : $94{,}2:4\times3=\res{70{,}65\ \mathrm{m}}$ environ.
\label{fig-cor-10-14}\figureGeo[corrige]{lon-14}
\end{corrige}\end{blocLivre}
\begin{blocLivre}{cor-10-15}\begin{corrige}{15}{Vrai ou faux}\demarche{Chercher un contre-exemple.}
1. $\res{\text{Vrai}}$ : $4\times(2\times c)=2\times(4\times c)$. 2. $\res{\text{Faux}}$ : un rectangle de 5 cm sur 1 cm et un carré de 3 cm de côté ont tous deux un périmètre de 12 cm.
\end{corrige}\end{blocLivre}
